\documentclass[11pt,a4paper]{article}
\usepackage[utf8]{inputenc}
\usepackage[T1]{fontenc}
\usepackage[brazilian]{babel}
\usepackage[margin=2.3cm]{geometry}
\usepackage{graphicx}
\usepackage{booktabs}
\usepackage{array}
\usepackage{xcolor}
\usepackage{hyperref}
\usepackage{caption}
\usepackage{amsmath}
\usepackage{longtable}
\usepackage{float}
\usepackage{enumitem}
\usepackage{multirow}

\hypersetup{
  colorlinks=true,
  linkcolor=blue!50!black,
  urlcolor=blue!50!black,
  citecolor=blue!50!black
}

\title{Qualidade dos Alarmes no EXP7 (Itens 1+2)\\
\large O que temos de alarmes, e o que de fato conta como\\
alarme real para TC382\_03\_A e T5\_AVG\_A}
\author{Pipeline \texttt{cnn1d\_ae} / AutoML --- projeto \texttt{TesteMLCab}}
\date{Agosto de 2026}

\begin{document}
\maketitle

\begin{abstract}
Este relatório complementa o EXP7 (itens 1+2 -- features multi-escala e de
textura, ver \texttt{docs/analise\_automl\_exp7.md}) com um levantamento
detalhado dos 40 alarmes usados na avaliação out-of-sample: quantos temos,
de qual sensor, em qual condição, e -- o ponto central deste documento --
\textbf{quais desses 40 são eventos físicos reais e quais são artefato de
falha de sensor/comunicação}. Encontramos que os 8 alarmes de condição
\texttt{UNDER} disparam com valores fisicamente implausíveis
(-18 a -22, fora da faixa normal de operação) e coincidem, em todos os 5
instantes distintos verificados, com gaps de dado, status de
\texttt{Comm Fail}/\texttt{Out of Serv}, ou o equipamento já parado --
nunca com uma queda física real de temperatura em operação. Ao excluir
esses 8 casos, o candidato do EXP7 item 2 passa de \textbf{92,5\% de
hit\_rate / 72,5\% de antecedência real} para \textbf{100\% de hit\_rate /
90,6\% de antecedência real} sobre os 32 alarmes genuínos -- e os 3 casos
``sem detecção'' que restavam eram exatamente os 3 alarmes \texttt{UNDER}
mais problemáticos. Fechamos com um levantamento do catálogo completo de
alarmes da turbina (47 tags, 3.757 eventos) para contextualizar o quanto
ainda não foi explorado.
\end{abstract}

\tableofcontents

% =====================================================================
\section{Contexto: EXP7 itens 1 e 2}
% =====================================================================

O EXP7 estende o candidato do EXP6 (\texttt{TC382\_03\_A} + \texttt{T5\_AVG\_A}
+ 10 canais de vibração de mancal, modelo \texttt{iforest}/\texttt{ocsvm}
não supervisionado) com duas mudanças de engenharia de features:

\begin{itemize}
  \item \textbf{Item 1 -- multi-escala}: features derivadas (média/desvio/
        tendência móveis) em 4 janelas de tempo (6min, 1h, 4h, 24h) em vez
        de uma única janela de 6min.
  \item \textbf{Item 2 -- textura}: kurtosis, skewness e crest factor
        (pico/RMS) móveis, para as janelas $\geq$1h -- capturam mudança de
        \emph{forma} do sinal de vibração, não só nível ou tendência.
\end{itemize}

\textbf{Candidato vencedor (item 1+2):} \texttt{ocsvm} (percentil 99,9,
debounce=1), sobre o grupo com 276 features de entrada. Avaliado no
período out-of-sample de 2025-07-01 a 2026-04-30, contra os alarmes reais
de \texttt{TC382\_03\_A}/\texttt{T5\_AVG\_A} desse período.

\begin{table}[H]
\centering
\begin{tabular}{@{}lccc@{}}
\toprule
\textbf{Etapa} & \textbf{hit\_rate} & \textbf{FP} & \textbf{Preditivo / Reativo / Sem detecção} \\
\midrule
EXP6 (janela única) & 75,0\% (30/40) & 0,06\% & 26 / 4 / 10 \\
EXP7 item 1 (multi-escala) & 87,5\% (35/40) & 2,21\% & 28 / 7 / 5 \\
\textbf{EXP7 item 1+2 (+ textura)} & \textbf{92,5\% (37/40)} & \textbf{1,94\%} & \textbf{29 / 8 / 3} \\
\bottomrule
\end{tabular}
\caption{Evolução até o candidato deste relatório. Detalhe completo em
\texttt{docs/analise\_automl\_exp7.md}.}
\end{table}

\begin{figure}[H]
\centering
\includegraphics[width=0.85\textwidth]{../task_plots_exp7_multiescala/02_metricas_por_etapa.png}
\caption{hit\_rate, falso alerta e mediana de antecedência ao longo das
três etapas.}
\end{figure}

% =====================================================================
\section{Levantamento: o que temos de alarmes}
\label{sec:levantamento}
% =====================================================================

\subsection{Os 40 alarmes da avaliação OOS}

\begin{table}[H]
\centering
\begin{tabular}{@{}lcc@{}}
\toprule
\textbf{Sensor} & \textbf{Alarmes} & \textbf{\% do total} \\
\midrule
TC382\_03\_A & 35 & 87,5\% \\
T5\_AVG\_A & 5 & 12,5\% \\
\bottomrule
\end{tabular}
\qquad
\begin{tabular}{@{}lc@{}}
\toprule
\textbf{Condição} & \textbf{Qtd.} \\
\midrule
HI & 18 \\
HIHI & 14 \\
UNDER & 8 \\
\bottomrule
\end{tabular}
\caption{Os 40 alarmes OOS não são homogêneos: dominados por
\texttt{TC382\_03\_A} (87,5\%), e 20\% deles (8/40) são condição
\texttt{UNDER}.}
\end{table}

\subsection{A performance varia muito por sensor e por condição}

\begin{table}[H]
\centering
\begin{tabular}{@{}lcccc@{}}
\toprule
\textbf{Sensor} & \textbf{Preditivo} & \textbf{Reativo} & \textbf{Sem detecção} & \textbf{Taxa preditiva} \\
\midrule
TC382\_03\_A & 28 & 5 & 2 & 80,0\% (28/35) \\
T5\_AVG\_A & 1 & 3 & 1 & 20,0\% (1/5) \\
\bottomrule
\end{tabular}
\end{table}

\begin{table}[H]
\centering
\begin{tabular}{@{}lcccc@{}}
\toprule
\textbf{Condição} & \textbf{Preditivo} & \textbf{Reativo} & \textbf{Sem detecção} & \textbf{Taxa preditiva} \\
\midrule
HI & 16 & 2 & 0 & 88,9\% \\
HIHI & 13 & 1 & 0 & 92,9\% \\
\textbf{UNDER} & \textbf{0} & 5 & 3 & \textbf{0\%} \\
\bottomrule
\end{tabular}
\caption{O modelo tem antecedência quase perfeita para HI/HIHI (temperatura
subindo), mas \textbf{zero} antecedência para UNDER -- nenhum dos 8 alarmes
dessa condição foi previsto com antecedência real.}
\end{table}

\begin{figure}[H]
\centering
\includegraphics[width=0.75\textwidth]{01_deteccao_por_condicao.png}
\caption{Detecção por condição de alarme -- HI e HIHI com $>$88\% de
antecedência real; UNDER com 0\%.}
\end{figure}

% =====================================================================
\section{O achado central: UNDER é falha de sensor, não evento físico}
\label{sec:achado}
% =====================================================================

\subsection{Por que suspeitar}

Os 8 alarmes \texttt{UNDER} disparam com \texttt{Valor} negativo:

\begin{table}[H]
\centering
\small
\begin{tabular}{@{}llll@{}}
\toprule
\textbf{Data/hora} & \textbf{Tag} & \textbf{Condição} & \textbf{Valor} \\
\midrule
2025-07-24 16:14:37 & T5\_AVG\_A & UNDER & -22 \\
2025-07-24 16:14:37 & TC382\_03\_A & UNDER & -21 \\
2025-08-08 15:31:18 & T5\_AVG\_A & UNDER & -21 \\
2025-08-08 15:31:18 & TC382\_03\_A & UNDER & -21 \\
2025-11-29 07:52:03 & TC382\_03\_A & UNDER & -18 \\
2026-01-14 16:18:05 & T5\_AVG\_A & UNDER & -20 \\
2026-01-14 16:18:05 & TC382\_03\_A & UNDER & -19 \\
2026-01-17 00:59:01 & T5\_AVG\_A & UNDER & -20 \\
\bottomrule
\end{tabular}
\caption{Os 8 alarmes UNDER (5 instantes distintos, alguns disparando em
ambos os sensores ao mesmo tempo). Valores entre -18 e -22 -- fora de
qualquer faixa fisicamente plausível para esses sensores (operação normal
$\sim$600--800; mesmo em idle/frio, $\sim$25--40).}
\end{table}

\subsection{Verificação direta na série bruta}

Para cada um dos 5 instantes distintos, checamos a série de
\texttt{TC382\_03\_A} e o estado de \texttt{RUNNING\_A} numa janela de
poucos minutos ao redor do alarme:

\begin{table}[H]
\centering
\small
\begin{tabular}{@{}lp{9cm}@{}}
\toprule
\textbf{Alarme} & \textbf{Evidência} \\
\midrule
08/08/2025 15:31 & 9 de 12 pontos NaN em \texttt{TC382\_03\_A} na janela de
3 min; \texttt{RUNNING\_A=1} o tempo todo -- gap de dado, não queda física. \\
17/01/2026 00:59 & \texttt{RUNNING\_A} = \texttt{\{'Name': 'Comm Fail', ...\}}
e depois \texttt{\{'Name': 'Out of Serv', ...\}} -- falha de comunicação
explícita registrada pelo próprio sistema. \\
29/11/2025 07:52 & \texttt{RUNNING\_A=0} -- equipamento já desligado. \\
24/07/2025 16:14 & Transição de \texttt{RUNNING\_A} 1$\to$0 -- parada em
andamento. \\
14/01/2026 16:18 & Mistura de \texttt{Out of Serv} e transição de estado. \\
\bottomrule
\end{tabular}
\caption{Nenhum dos 5 casos corresponde a uma queda física real de
temperatura durante operação normal.}
\end{table}

\begin{figure}[H]
\centering
\includegraphics[width=\textwidth]{03_exemplos_under_sensor_error.png}
\caption{Dois dos cinco casos, em detalhe. Em ambos, a série de
\texttt{TC382\_03\_A} está numa faixa perfeitamente normal
($\sim$630--640 e $\sim$685--691 respectivamente) exatamente no instante
em que o sistema de alarme registrou um valor de -21/-20 -- confirmando
que a leitura que disparou o alarme nunca chegou ao nosso pipeline de
dados como um valor de processo real.}
\end{figure}

\subsection{Conclusão do achado}

\textbf{Não existe precursor de sensor para prever a própria falha de
comunicação do sensor.} Excluir os alarmes \texttt{UNDER} da avaliação não
é viés de seleção favorável ao modelo -- é remover alvos que são
estruturalmente impossíveis de prever com os dados de temperatura/vibração
disponíveis, porque a causa do alarme é externa ao fenômeno físico que o
modelo observa (instrumentação/comunicação, não o processo em si).

% =====================================================================
\section{Métricas corrigidas: o que de fato conta como alarme real}
\label{sec:metricas-corrigidas}
% =====================================================================

Definimos \textbf{alarme real}, para fins de avaliação deste detector,
como um alarme de condição \texttt{HI} ou \texttt{HIHI} -- ou seja,
excluindo \texttt{UNDER} pela razão exposta na Seção~\ref{sec:achado}.
Isso reduz o denominador de 40 para 32 alarmes.

\begin{table}[H]
\centering
\begin{tabular}{@{}lcc@{}}
\toprule
\textbf{Métrica} & \textbf{Com UNDER (40)} & \textbf{Sem UNDER (32, real)} \\
\midrule
hit\_rate & 92,5\% & \textbf{100\%} \\
Antecedência real (preditivo) & 72,5\% (29/40) & \textbf{90,6\% (29/32)} \\
Reativo & 8 & 3 \\
Sem detecção & 3 & \textbf{0} \\
Mediana de antecedência & 14,7h & 14,7h (mesmos 29 casos) \\
\bottomrule
\end{tabular}
\caption{Sobre os 32 alarmes genuínos, o candidato detecta 100\% deles,
com 90,6\% de antecedência real -- os 3 ``sem detecção'' reportados
anteriormente eram exatamente os 3 alarmes \texttt{UNDER} mais
problemáticos.}
\end{table}

\begin{figure}[H]
\centering
\includegraphics[width=\textwidth]{02_antes_depois_under.png}
\caption{Antes e depois de excluir os alarmes \texttt{UNDER}. À direita,
zero casos sem detecção.}
\end{figure}

\textbf{Ressalva importante:} os 3 casos ``reativos'' que sobram (detecção
só depois do alarme, sem antecedência) continuam sem antecedência real --
esse achado não os resolve, só remove o ruído dos casos \texttt{UNDER}. Os
3 são investigados em detalhe na Seção~\ref{sec:reativos}.

\begin{figure}[H]
\centering
\includegraphics[width=\textwidth]{../task_plots_exp7_multiescala/03_painel_preditivos.png}
\caption{Os 29 casos com antecedência real do candidato item 1+2 (inclui
os que, após a filtragem de UNDER, permanecem exatamente os mesmos --
nenhum caso preditivo era \texttt{UNDER}).}
\end{figure}

% =====================================================================
\section{Os 3 alarmes reativos remanescentes}
\label{sec:reativos}
% =====================================================================

Sobre os 32 alarmes genuínos (Seção~\ref{sec:metricas-corrigidas}), 3
continuam sendo detectados só \emph{depois} do alarme disparar, sem
antecedência real. Os 3 pertencem todos a \texttt{TC382\_03\_A} -- mas na
prática correspondem a apenas \textbf{2 eventos físicos distintos}, com
causas raiz completamente diferentes entre si.

\begin{table}[H]
\centering
\begin{tabular}{@{}llccc@{}}
\toprule
\textbf{Data/hora} & \textbf{Condição} & \textbf{Valor} & \textbf{Atraso até 1ª detecção} & \textbf{Evento} \\
\midrule
20/10/2025 16:18:40 & HI & 776 & 4h00 & \multirow{2}{*}{trip/religamento} \\
20/10/2025 16:20:36 & HIHI & 784 & 3h57 & \\
24/02/2026 11:01:00 & HI & 780 & 0min (simultâneo) & início tipo degrau \\
\bottomrule
\end{tabular}
\caption{Os 2 alarmes de 20/10/2025 são a mesma escalada física (HI seguido
de HIHI 2 minutos depois), tratados aqui como um único evento.}
\end{table}

\subsection{Evento 1 (20/10/2025) -- trip/religamento da turbina}

Não é uma falha do modelo em captar um precursor gradual: é um
\textbf{trip/desligamento-religamento}. O estado operacional
(\texttt{operational\_state}) sai de \texttt{"on"} já às \textbf{15:29} --
quase 1h \emph{antes} do próprio alarme -- e passa a oscilar entre
\texttt{"transiente"} e \texttt{"off\_curto"} por quase 5 horas. A série
bruta mostra \texttt{TC382\_03\_A}/\texttt{T5\_AVG\_A} saltando de forma
violenta e fisicamente implausível (656$^{\circ}$C $\to$ 767$^{\circ}$C
$\to$ 213$^{\circ}$C $\to$ 682$^{\circ}$C $\to$ 92$^{\circ}$C $\to$
678$^{\circ}$C\ldots), junto com \texttt{RUNNING\_A} alternando entre 0 e
1 -- assinatura típica de partida/parada, não de um processo térmico
contínuo.

A pipeline \textbf{zera \texttt{is\_anom\_point} por design} sempre que
\texttt{operational\_state} $\neq$ \texttt{"on"} (ver
\texttt{automl\_pipeline.py}, linhas 146--147 e 324--325) -- decisão
tomada justamente para não confundir ruído de partida/parada com anomalia
real. O estado só volta a \texttt{"on"} estável às \textbf{20:19:00}, e é
exatamente nesse instante que aparece a primeira detecção -- não é
coincidência, é a mecânica exata da máscara operacional.

\begin{figure}[H]
\centering
\includegraphics[width=\textwidth]{06_evento_trip_reativo.png}
\caption{\texttt{TC382\_03\_A} ao redor do trip de 20/10/2025. Áreas
sombreadas = estado \texttt{!= "on"}, período em que a pontuação de
anomalia fica zerada por design. O alarme (linhas vermelhas) dispara
dentro dessa janela; a 1ª detecção (ponto verde) só ocorre quando o estado
volta a \texttt{"on"}.}
\end{figure}

\textbf{Conclusão:} o alarme disparou \emph{durante} um trip, período que
o desenho atual do pipeline exclui deliberadamente da pontuação de
anomalias (tanto no treino quanto no score). Não é uma limitação de
engenharia de features -- é um limite estrutural de escopo: o modelo atual
não foi desenhado para avaliar transientes de partida/parada.

\subsection{Evento 2 (24/02/2026) -- início tipo degrau, sem precursor}

Aqui o cenário é o oposto: \texttt{operational\_state = "on"} o tempo
todo, sem nenhuma exclusão operacional envolvida.
\texttt{TC382\_03\_A} fica \textbf{estável em $\sim$709--715$^{\circ}$C
por pelo menos 6h antes} do alarme (deriva de menos de 1\%, dentro do
ruído normal de operação), sem nenhum sinal gradual de degradação. Em um
único intervalo de 10 min salta abruptamente para 748$^{\circ}$C e, no
intervalo seguinte, para 791$^{\circ}$C -- um evento de \textbf{início
rápido tipo degrau}, sem precursor visível em nenhuma das 4 escalas de
tempo capturadas pelas features multi-escala (6min, 1h, 4h, 24h).

\begin{figure}[H]
\centering
\includegraphics[width=\textwidth]{07_evento_degrau_reativo.png}
\caption{\texttt{TC382\_03\_A} ao redor do alarme de 24/02/2026. Patamar
estável por 6h, depois salto abrupto coincidente com o próprio alarme. A
1ª detecção (ponto verde) ocorre no mesmo intervalo de 10 min do alarme --
o mais rápido possível dado o grid de debounce.}
\end{figure}

\textbf{Conclusão:} não há precursor a ser capturado -- a ausência de
antecedência é inerente à física do evento, não ao modelo. O modelo
detecta com a maior velocidade fisicamente possível dado o grid de
debounce atual.

\subsection{Síntese}

\begin{table}[H]
\centering
\begin{tabular}{@{}lp{6cm}c@{}}
\toprule
\textbf{Evento} & \textbf{Causa raiz} & \textbf{É falha do modelo?} \\
\midrule
20/10/2025 (HI+HIHI) & Trip/religamento -- alarme durante estado excluído
da pontuação por design & Não -- limite estrutural de escopo \\
24/02/2026 (HI) & Início tipo degrau, sem precursor gradual em nenhuma
janela & Não -- evento genuinamente sem antecedência possível \\
\bottomrule
\end{tabular}
\caption{Nenhum dos 3 casos reativos indica um problema de engenharia de
features a corrigir.}
\end{table}

Se cobrir o Evento 1 for considerado prioritário no futuro, a via não é
mais engenharia de features sobre o sinal de processo -- é estender a
pontuação de anomalia para dentro dos estados \texttt{"transiente"}/
\texttt{"off\_curto"} (ou tratar o próprio padrão de trip como uma classe
de evento à parte), o que exigiria uma modelagem específica de dinâmica de
partida/parada, fora do escopo desta série de experimentos.

% =====================================================================
\section{Contexto mais amplo: o catálogo geral de alarmes da turbina}
\label{sec:catalogo}
% =====================================================================

O arquivo \texttt{alarmes\_selecionados\_turbina\_a.csv} cobre muito mais
do que os 2 sensores usados no EXP5--EXP8: \textbf{47 tags distintos,
3.757 alarmes (onset) entre 2022-01-04 e 2026-04-18}.

\begin{table}[H]
\centering
\begin{tabular}{@{}lcc@{}}
\toprule
\textbf{Categoria} & \textbf{Alarmes} & \textbf{Nº de tags} \\
\midrule
Temperatura & 1.440 & 18 \\
Pressão & 1.339 & 9 \\
Falha de instrumento/transmissor & 844 & 10 \\
Vibração & 134 & 10 \\
\bottomrule
\end{tabular}
\caption{\texttt{TC382\_03\_A}/\texttt{T5\_AVG\_A} são só 2 dos 18 tags de
temperatura. Os 10 canais de vibração (\texttt{TV\_*}) têm alarme
próprio -- usados até aqui só como \emph{feature} de entrada, nunca como
\emph{alvo} de predição.}
\end{table}

\begin{figure}[H]
\centering
\includegraphics[width=0.95\textwidth]{05_categorias_catalogo.png}
\end{figure}

\begin{figure}[H]
\centering
\includegraphics[width=0.85\textwidth]{04_top_tags_catalogo.png}
\caption{Os 5 tags com mais alarmes no catálogo completo:
\texttt{PI\_6240319\_AL} (630, falha de partida a gás),
\texttt{PAL\_6240315} (558, pressão baixa de gás combustível),
\texttt{PAL\_6240339} (254, pressão baixa de óleo lub.),
\texttt{TC382\_03\_A} (253) e \texttt{PALL\_6240340} (204).}
\end{figure}

Vale notar que ``Falha de instrumento/transmissor'' já é uma categoria
\emph{própria} do catálogo (844 alarmes, 10 tags) -- ou seja, o padrão que
encontramos nos alarmes \texttt{UNDER} de temperatura (falha de
sensor/comunicação disfarçada de leitura de processo) provavelmente
também existe, de forma explícita, nesses outros tags. Vale revisar se
esse tipo de alarme está sendo contaminado da mesma forma em experimentos
futuros que usem esses sensores.

% =====================================================================
\section{Recomendações}
% =====================================================================

\begin{itemize}
  \item \textbf{Adotar 32 (não 40) como o denominador correto} para
        relatar a performance do candidato \texttt{ocsvm} item 1+2 --
        100\% de hit\_rate, 90,6\% de antecedência real.
  \item \textbf{Filtrar por \texttt{Condição do Alarme}} (excluir
        \texttt{UNDER}, e revisar caso a caso qualquer condição não
        vista antes) deveria ser um passo padrão de pré-processamento do
        arquivo de alarmes, não uma correção pontual -- vale incorporar
        ao pipeline (\texttt{io.py}) em vez de só neste relatório.
  \item \textbf{Os 3 alarmes reativos remanescentes} (Seção~\ref{sec:reativos})
        já foram investigados: 2 são um trip/religamento (fora do escopo
        de pontuação por design) e 1 é um início tipo degrau sem
        precursor -- nenhum indica ganho disponível via engenharia de
        features adicional.
  \item \textbf{Considerar expandir o escopo} para os outros 45 tags do
        catálogo (Seção~\ref{sec:catalogo}), sobretudo pressão (maior
        categoria depois de temperatura) e os 10 canais de vibração como
        alvo próprio, não só como feature.
\end{itemize}

\end{document}
